\section{Introduction}

The \textit{Disquisitiones Arithmeticae} (1801) is Gauss's magnum opus, containing 355 articles across seven sections. It founded modern number theory and introduced the systematic use of congruences. In this chapter, we explore selected topics from this masterpiece, implementing Gauss's algorithms and extending his ideas with modern computation.

Gauss began writing the \textit{Disquisitiones} in 1795 at the age of 17 and published it in 1801. The work was so influential that it earned Gauss the epithet ``Princeps Mathematicorum'' (Prince of Mathematicians). The seven sections cover: congruences, linear congruences, quadratic residues, higher congruences, cyclotomy, quadratic forms, and application of cyclotomy to geometry.

\section{Section I: Congruences (Articles 1--45)}

We covered the basics in Chapter~\ref{ch:modular}. Here we highlight:
\begin{itemize}
    \item Article 36--44: Chinese Remainder Theorem (Gauss's proof)
    \item Article 45: Solution of general linear congruences $ax \equiv b \pmod{n}$
\end{itemize}

\section{Section II: Linear Congruences (Articles 46--59)}

\begin{theorem}[Article 49]
The congruence $ax \equiv b \pmod{n}$ has a solution iff $\gcd(a,n) \mid b$. If so, there are exactly $\gcd(a,n)$ solutions modulo $n$.
\end{theorem}

\section{Section III: Quadratic Residues (Articles 107--149)}

This section culminates in the Law of Quadratic Reciprocity (Article 131). Key articles:
\begin{itemize}
    \item Article 108--111: Definition and basic properties of Legendre symbol
    \item Article 129: Gauss's Lemma
    \item Article 131: Quadratic Reciprocity Law
    \item Article 134: Supplementary laws for $(-1/p)$ and $(2/p)$
    \item Article 142--149: Applications and examples
\end{itemize}

\section{Section IV: Gaussian Sums (Articles 356--358)}

Gauss's fourth proof of quadratic reciprocity using Gauss sums.

\section{Section V: Cyclotomy (Articles 355--366)}

Gauss's investigation of the division of the circle, i.e., roots of unity \index{roots of unity}. This leads to:

\begin{theorem}[Gauss's Criterion for Constructible Polygons]
A regular $n$-gon is constructible with ruler and compass iff $n = 2^k p_1 p_2 \cdots p_m$ where each $p_i$ is a distinct Fermat prime \index{Fermat prime} ($3, 5, 17, 257, 65537, \ldots$).
\end{theorem}

\section{Section VI: Quadratic Forms (Articles 150--300)}

This is the heart of the \textit{Disquisitiones}. Gauss developed the theory of binary quadratic forms $ax^2 + bxy + cy^2$.

\begin{definition}[Binary Quadratic Form \index{binary quadratic form}]
A binary quadratic form is $f(x,y) = ax^2 + bxy + cy^2$ with $a,b,c \in \Z$. Its \textit{discriminant} \index{discriminant!quadratic form} is $\Delta = b^2 - 4ac$.
\end{definition}

\begin{theorem}[Gauss's Composition Law \index{composition law!quadratic forms}]
The set of equivalence classes of primitive binary quadratic forms of fixed discriminant $\Delta < 0$ forms a finite abelian group under composition.
\end{theorem}

This group is isomorphic to the ideal class group \index{ideal class group} \index{class group!quadratic forms} of the quadratic field $\Q(\sqrt{\Delta})$.

\begin{theorem}[Dirichlet's Unit Theorem]\index{Dirichlet's unit theorem}
The unit group of the ring of integers of a quadratic field is of the form $\mu \times \Z$, where $\mu$ is the finite cyclic group of roots of unity.
\end{theorem}

\begin{theorem}[Minkowski's Bound]\index{Minkowski bound}
Every ideal class of a quadratic field contains an ideal with norm bounded by $\frac{2}{\pi}\sqrt{|\Delta|}$.
\end{theorem}

\section{Section VII: Cyclotomic Equations (Articles 355--366)}

Gauss solved $x^n - 1 = 0$ by radicals, showing the constructibility of the regular 17-gon (Article 365).

\section{Python Implementation}

The implementation in \texttt{python/disquisitiones.py} includes:

\begin{lstlisting}[style=python, caption={Key functions in python/disquisitiones.py}]
solve_linear_congruence(a, b, n)      # Article 49
quadratic_forms(discriminant)         # Generate forms of given disc
compose_forms(f, g)                   # Gauss composition
class_group(discriminant)             # Class group structure
cyclotomic \index{cyclotomy}polynomial(n)              # Phi_n(x)
heptadecagon_construction()           # Gauss's 17-gon
euler_criterion(a, p)                 # Article 132
legendre_symbol(a, p)                 # Quadratic residue test
gauss_lemma(a, p)                     # Article 129
\end{lstlisting}

\section{Numerical Examples}

\begin{example}[Binary Quadratic Forms]
Discriminant $\Delta = -20$. Reduced forms:
\begin{lstlisting}[style=python]
>>> from disquisitiones import reduced_forms
>>> reduced_forms(-20)
[(1, 0, 5), (2, 2, 3)]  # x^2 + 5y^2, 2x^2 + 2xy + 3y^2
\end{lstlisting}
Class number $h(-20) = 2$.
\end{example}

\begin{example}[Class Group Computation]
\begin{lstlisting}[style=python]
>>> from disquisitiones import class_group
>>> class_group(-23)
Cyclic group of order 3 generated by form (2, 1, 3)
\end{lstlisting}
\end{example}

\section{Exercises}

\begin{exercise}
Implement Gauss's composition algorithm for binary quadratic forms (Article 234--238). Verify it gives a group structure.
\end{exercise}

\begin{exercise}
Compute the class number \index{class number} $h(\Delta)$ for all $\Delta > -100$. Plot the distribution.
\end{exercise}

\begin{exercise}
Implement Gauss's algorithm for the 17-gon construction. Output coordinates of vertices.
\end{exercise}

\begin{exercise}
Prove that the class group of discriminant $\Delta$ is isomorphic to the ideal class group of $\Q(\sqrt{\Delta})$.
\end{exercise}

\begin{exercise}
Verify Dirichlet's unit theorem for $\Q(\sqrt{2})$ and $\Q(\sqrt{-1})$. Find the fundamental units.
\end{exercise}



\section{Exercise Solutions}

\begin{solution}
Gauss's composition algorithm (Articles 234--238) takes two forms $(a,b,c)$ and $(a',b',c')$ with the same discriminant and produces a third form representing the product of ideals. The algorithm finds $B \equiv b \pmod{2a}$, $B \equiv b' \pmod{2a'}$, $B^2 \equiv \Delta \pmod{4aa'}$, then sets the composed form to $(aa'/g, B, \ldots)$ where $g = \gcd(a, a', (B+\Delta)/2)$. Verification shows the composition is associative with identity and inverses.
\end{solution}

\begin{solution}
Computing $h(\Delta)$ for $-100 < \Delta < 0$ using reduced forms: the class numbers are small (mostly 1--4) with occasional spikes at discriminants with many prime factors. The distribution shows that most fundamental discriminants have class number 1 or 2. The plot reveals no simple pattern, consistent with the class number formula $h(\Delta) = \frac{w\sqrt{|\Delta|}}{2\pi} L(1, \chi_\Delta)$.
\end{solution}

\begin{solution}
The 17-gon construction uses the period equations of the 17th cyclotomic field. Gauss showed that $\zeta + \zeta^{-1} = 2\cos(2\pi/17)$ satisfies a quadratic over $\Q(\zeta+\zeta^{16}, \zeta^4+\zeta^{13}, \ldots)$. Iterating this process gives nested square roots. The coordinates are $(\cos(2\pi k/17), \sin(2\pi k/17))$ for $k = 0, \ldots, 16$.
\end{solution}

\begin{solution}
The map from ideal classes to form classes sends an ideal $\mathfrak{a} = [a, \frac{b+\sqrt{\Delta}}{2}]$ to the form $a x^2 + bxy + \frac{b^2-\Delta}{4a} y^2$. This is well-defined, surjective, and injective (by reduction theory). The group operation corresponds to ideal multiplication, establishing the isomorphism $\operatorname{Cl}(\Q(\sqrt{\Delta})) \cong \operatorname{Cls}(\Delta)$.
\end{solution}

\begin{solution}
For $\Q(\sqrt{2})$, the fundamental unit is $1+\sqrt{2}$ since $(1+\sqrt{2})(-1+\sqrt{2}) = 1$. For $\Q(\sqrt{-1}) = \Q(i)$, the unit group is $\{\pm 1, \pm i\} \cong \Z/4\Z$, consistent with $\mu = 4$ roots of unity. Dirichlet's theorem states $\mathcal{O}_K^* \cong \mu \times \Z^{r_1+r_2-1}$ where $r_1$ is the number of real embeddings and $r_2$ the number of complex pairs.
\end{solution}


\section{References}

\begin{itemize}
    \item \cite{gauss-disq} -- Complete \textit{Disquisitiones Arithmeticae}
    \item \cite{cox-primes} -- Cox, \textit{Primes of the Form $x^2+ny^2$}
    \item \cite{buell-binary} -- Buell, \textit{Binary Quadratic Forms}
    \item \cite{cohen-forms} -- Cohen, \textit{Number Theory}, Vol. 2
    \item \cite{neukirch} -- Neukirch, \textit{Algebraic Number Theory}
\end{itemize}